\subsection{P002 --- Graph-based spatial–temporal prediction and feature interaction analysis of CO2 and occupant in large indoor space}
\label{paper:P002}
\textbf{Authors:} Cong Huang, Helen H.L. Kwok, Kwok Ho Poon, Zhaoji Wu, Fangli Hou, Jun Ma, Jack C.P. Cheng\par
\textbf{Major Category:} Building Environment and IEQ\par
\textbf{Secondary Category:} AI and Machine Learning for Buildings, Occupancy and Human-Building Interaction, Sensors, IoT and Data Integration\par
\textbf{Keywords:} Graph neural network, Spatial-temporal prediction, Spatial-temporal feature interactions, Indoor Air Quality, Indoor occupant, Dynamic graph learning, Occupant-in-loop\par
\paragraph{主な主張}
提案AST-GCGRUは, 6 hおよび12 hのmulti-step predictionでgraph-level R2 0.898および0.812を示し, GCGRUの0.860および0.699を上回った. 12 hではoccupant node R2 0.59, CO2 node R2 0.73で, GCGRUの0.32, 0.58より高かった. \cite{Huang:2025GraphCO2Occupant}
\paragraph{新規性}
距離ベースの固定graphではなく, partial correlationで初期化したgraphをST-Attで学習中に更新し, time-lag nodesにより30 min CO2 sensorと6 min occupant sensorの異なる時間解像度を保持したまま, CO2とoccupancyを同一graphで同時予測・interaction analysisする構成を提示した.
\paragraph{位置付け}
大空間におけるIAQとoccupancyの同時予測・相互作用分析を主目的とするため, Building Environment and IEQを主分類とし, AI and Machine Learning for Buildings, Occupancy and Human-Building Interaction, Sensors, IoT and Data Integrationを副分類として位置付ける.
